\chapter{绪论}

\section{研究背景与意义}

医学图像分割是计算机辅助诊断、术前规划与放疗靶区勾画中的关键基础能力\cite{litjens2017survey}。高质量的分割结果直接影响器官体积估计、病灶边界判读及下游治疗决策。相较于自然图像任务，医学分割长期面临两类共性难题：其一为训练信号分配问题，即目标区域相对背景高度稀疏、困难像素数量有限，致使模型训练易被大量简单像素主导；其二为局部结构表征问题，即边界模糊、低对比、细长或小体积结构往往决定临床可用性，但在以全局建模为主的网络中难以获得充分刻画\cite{hesamian2019survey}。

深度学习的发展持续推动医学分割方法的演进。从全卷积网络（FCN）\cite{long2015fcn}、U-Net\cite{ronneberger2015unet}到 nnU-Net\cite{isensee2021nnunet}，再到基于 Transformer 的 TransUNet\cite{chen2021transunet}与 Swin-Unet\cite{cao2022swinunet}，模型的表征能力与适配范围不断提升。近年来，Segment Anything Model（SAM）\cite{kirillov2023segment}提出了统一的提示式分割范式，其后续版本 SAM 2\cite{ravi2024sam2}进一步扩展至视频场景。MedSAM\cite{ma2024medsam}将 SAM 迁移至医学图像领域，显著增强了模型的通用提示分割能力。然而，当 MedSAM 应用于腹部多器官 CT 分割时，仍暴露出两方面关键不足：一是类别不平衡与困难像素学习不足并存，二是 ViT 编码器对局部边界与高频纹理的建模能力相对欠缺。围绕上述问题开展系统优化，兼具理论研究价值与实际工程意义。

从应用层面而言，腹部多器官分割是医学图像分析中最具代表性的复杂场景之一。其难点不仅在于类别数目较多，更在于不同器官在体积尺度、形态规则性、灰度对比度及邻接关系上存在显著差异。例如，肝脏、脾脏等大器官通常具有较稳定的解剖轮廓，而胰腺、食管、十二指肠及肾上腺等器官往往体积较小、边界模糊、形态不规则，更易成为模型误差的主要来源。因此，在该场景中验证方法，既能检验模型的整体分割能力，又能敏感地暴露训练信号设计与局部结构建模方面的不足。

从研究层面而言，视觉基础模型的迁移正在改变医学图像分割的技术路径。一方面，MedSAM 等模型已具备较强的通用表示与提示驱动能力，使研究者无需从零构建任务专用网络；另一方面，基础模型的强大并不意味着下游任务不再需要针对性优化。对于医学分割而言，如何在不破坏预训练表示优势的前提下进一步重构训练信号、补足局部边界表征、提升困难器官的可分性，仍是一个值得深入探讨的问题。本文正是在这一背景下展开研究：并非试图构建一种全新的医学分割框架，而是围绕 MedSAM 在具体任务中的能力边界进行系统性优化。

从论文工作的意义而言，本文的价值主要体现在两个方面。其一，在方法层面，本文分别从损失设计与结构补强两个方向出发，探讨 MedSAM 当前性能瓶颈的主要来源及优先介入方向；其二，在研究组织层面，本文并非将所有尝试简单堆叠为模块集合，而是通过统一实验协议下的逐步筛选，构建起一条由问题定义、路线比较、负向结果排除到最终方案保留的证据链。

\section{国内外研究现状}

\subsection{医学图像分割模型发展}

医学图像分割方法经历了从传统方法到深度学习，再到视觉基础模型迁移的多阶段演进\cite{wang2022medseg_survey}。

在深度学习时代之前，医学图像分割主要依赖于阈值、区域生长、活动轮廓模型等传统方法\cite{litjens2017survey}。这些方法虽然在特定场景下有效，但严重依赖手工特征设计和先验知识，面对复杂解剖结构和多变成像条件时泛化能力不足。

全卷积网络（FCN）\cite{long2015fcn}的提出标志着深度学习在语义分割中的突破。U-Net\cite{ronneberger2015unet}通过编码器—解码器结构和跳跃连接机制，奠定了医学分割的经典范式。此后，研究者围绕多尺度融合、注意力建模和三维体积信息利用等方向持续改进：UNet++\cite{zhou2018unetpp}通过密集跳跃路径缓解语义鸿沟；Attention U-Net\cite{oktay2018attention}引入注意力门控聚焦目标区域；3D U-Net\cite{cicek20163dunet}和 V-Net\cite{milletari2016vnet}进一步将体积上下文和 Dice 优化引入医学分割任务。

在自动化配置方面，nnU-Net\cite{isensee2021nnunet}通过预处理、网络拓扑和后处理的自适应配置成为医学分割的强基线。在多尺度建模方面，PSPNet\cite{zhao2017pspnet}、DeepLabV3\cite{chen2017deeplabv3}、DeepLabV3+\cite{chen2018deeplabv3plus}和 DenseNet\cite{huang2017densenet}等工作也为密集预测任务提供了重要技术积累。

随着 Transformer 架构进入视觉领域，医学分割进入新的阶段。代表性工作包括 TransUNet\cite{chen2021transunet}、SETR\cite{zheng2021setr}、Swin-Unet\cite{cao2022swinunet}、SegFormer\cite{xie2021segformer}和 UNETR\cite{hatamizadeh2022unetr}，这些方法共同验证了全局自注意力对长距离依赖建模的优势。此后，SAM\cite{kirillov2023segment}和 MedSAM\cite{ma2024medsam}推动了提示式分割在医学场景中的应用。SAM 2\cite{ravi2024sam2}进一步将分割能力扩展至视频和时序场景，但其医学领域适配仍处于早期探索阶段，目前已有的 SAM 2 医学适配工作主要集中在视频超声和内镜序列等时序场景上\cite{zhu2024medsam2}，尚未形成针对静态 CT 切片的成熟方案。围绕医学任务适配，SAMed\cite{zhang2024samed}、Medical SAM Adapter\cite{wu2023medsamadapter}、SAM-Adapter\cite{chen2023sam_survey_med}、SAM-Med2D\cite{cheng2023sammed2d}和 SAM-Med3D\cite{wang2023sammed3d}等工作进一步探索了基础模型的高效适配方式。已有研究普遍表明：基础模型迁移为医学分割提供了新起点，但任务特异性的损失设计与局部结构补强仍然必要\cite{mazurowski2023sam_survey}。

总体而言，医学图像分割模型的发展并非简单的新旧替代，而更体现为优化重心的持续迁移：早期研究主要解决目标区域的基本定位问题，随后转向多尺度结构与长距离依赖的稳健学习，而在基础模型出现后，研究重点进一步转变为在已有强初始化之上补足任务特异性短板。本文选择以 MedSAM 为基座而非从零设计独立网络，正是基于当前研究价值更集中于后一问题的判断。

\subsection{不平衡学习方法}

类别不平衡是医学图像分割中的长期核心问题\cite{litjens2017survey}。在腹部多器官分割中，背景像素通常占据绝大多数，而小器官和边界区域的有效监督信号相对稀缺。

针对类别间不平衡，Class-Balanced Loss\cite{cui2019classbalanced}通过有效样本数建模多数类信息冗余，给出了系统化的类别重加权方法。针对困难样本学习，OHEM\cite{shrivastava2016ohem}提出在线困难样本挖掘策略，Focal Loss\cite{lin2017focal}通过连续调制因子降低易样本权重，GHM\cite{li2019ghm}则从梯度密度角度重新审视样本加权问题。

在损失函数设计方面，Dice Loss\cite{milletari2016vnet}因其对不平衡数据的天然鲁棒性而被广泛使用。Tversky Loss\cite{salehi2017tversky}通过调节假阳性与假阴性权重提供更灵活的平衡机制，Boundary Loss\cite{kervadec2019boundary}则从距离度量角度直接优化边界质量，Active Boundary Loss\cite{wang2022activeboundaryloss}进一步通过动态聚焦边界区域提高边界预测精度。Lov\'{a}sz-Softmax\cite{berman2018lovasz}进一步提供了面向重叠指标优化的代理形式。Unified Focal Loss\cite{yeung2022unified}则尝试从统一建模角度整合不平衡学习策略。总体来看，现有方法大多从“重叠度优化”或“连续调制”角度缓解不平衡，但对类别间重平衡与类别内困难像素加权的联合、可解释建模仍有空间。本文第3章正是在这一问题上，对 MedSAM 的训练信号进行任务驱动的重构。

进一步看，现有不平衡学习方法虽然提供了丰富的技术工具，但其设计出发点往往并非完全针对提示式医学分割场景。一方面，许多方法默认在多类别语义分割或检测框架中使用，其多数类与少数类的划分更多基于类别频次；另一方面，本文任务是在 box prompt 约束下进行二值目标恢复，真正影响性能的因素不仅包括前景与背景的数量差异，还包括同一目标内部困难像素与易分像素之间的梯度竞争。因此，本文的目标并非简单引入某种更强的通用不平衡损失，而是将当前任务中最关键的双层失衡予以显式建模。

\subsection{参数高效微调与局部适配研究}

随着视觉基础模型参数规模持续增长，全参数微调的成本问题日益突出。参数高效微调（Parameter-Efficient Fine-Tuning, PEFT）技术通过冻结大部分预训练参数、仅更新少量任务相关参数来实现任务适配\cite{ding2022peftsurvey}。LoRA\cite{hu2022lora}通过低秩矩阵分解建模权重增量，在自然语言处理任务上取得了显著效果；Visual Prompt Tuning（VPT）\cite{jia2022vpt}、AdaptFormer\cite{chen2022adaptformer}和 Adapter\cite{houlsby2019adapter}等方法也从不同角度验证了轻量级适配的可行性。

然而，在医学图像分割这类强依赖空间重建与边界细节的任务中，单纯冻结主干并不总能维持性能。已有研究指出，ViT 的全局注意力与 CNN 的局部归纳偏置具有较强互补性\cite{dai2021coatnet}。在医学 SAM 适配方向，I-MedSAM\cite{wei2024imedsam}通过高频增强补充边界信息，MA-SAM\cite{chen2024masam}和 Medical SAM Adapter\cite{wu2023medsamadapter}则通过插入卷积适配模块提升空间建模能力。这些工作说明：在医学场景中，合理的结构补强可能比单纯追求参数效率更重要。本文第4章将围绕这一问题，对多条微调与适配路线进行系统比较，并最终收敛到局部适配方案。

这也意味着，对 PEFT 路线的评价标准不应仅停留在新增参数是否最少这一维度。对于需要高质量边界恢复的器官分割任务而言，若参数虽然节省但主干表示无法继续适应具体医学纹理和局部几何特征，则方法的综合收益反而可能受限。正因如此，本文在第4章中并不预设 LoRA 一定优于全参微调，亦不将局部适配器直接视为先验正确方案，而是将其统一纳入比较框架，以突出路线筛选本身的研究意义。

\subsection{视觉基础模型迁移带来的新变化}

视觉基础模型进入医学图像分析后，研究范式出现了明显变化。传统医学分割研究往往从网络结构本身出发，围绕编码器、解码器、跳跃连接或注意力模块进行任务专用设计；而以 SAM、MedSAM 为代表的基础模型路径，则更强调“先获得强通用表征，再进行任务适配”。这意味着研究重点不再只是设计一个从头训练的新网络，而是需要回答另一个问题：对于已经具备较强初始化能力的基础模型，下游性能瓶颈究竟还主要来自哪里。

这一变化带来了两个直接影响。其一，传统方法中许多从零建模的问题在基础模型场景下被部分弱化。例如，MedSAM 已能在较大范围内提供可用的分割初始化，因此简单重复构造新的主干网络未必能带来显著收益。其二，新的问题开始凸显，即如何在保留基础模型已有优势的同时，对医学任务的局部困难区域进行低破坏性补强。本文第3章与第4章的两条主线本质上均建立在这一认识之上：第3章回答基础模型在当前任务中尚缺何种训练信号，第4章回答在训练信号重构后尚缺何种局部表征。

因此，本文并非将基础模型视为需要被完全替换的对象，而更倾向于将其视为一个强基座：它已解决了诸多底层表征问题，但在特定任务中仍保留了若干可被继续优化的薄弱环节。围绕这些薄弱环节开展针对性改进，相较于从零设计复杂框架，更符合当前基础模型时代医学图像分割研究的现实路径。

本文不试图将结论直接外推为通用医学分割方法，而是以腹部多器官 CT 为统一验证场景，聚焦 MedSAM 在该类任务中的两个关键问题：

\begin{enumerate}
  \item \textbf{训练信号分配不足}：在 box prompt 驱动的二值分割范式下，前景像素相对背景高度稀疏，边界与低对比区域中的困难像素又占比较低，常规损失容易被大量简单像素主导。
  \item \textbf{局部边界表征不足}：ViT 编码器更擅长建模全局语义关系，但对边界、高频纹理和细粒度局部结构的响应相对不足，限制了复杂器官边界的进一步提升。
\end{enumerate}

围绕这两个问题，本文将研究目标定位为：以 MedSAM 为基座，构建一条”问题发现—针对性修正—逐步筛选—最终验证”的优化链。具体而言，先通过损失策略重构训练信号，再通过微调/适配路线比较补足局部边界表征，最终形成用于统一展示的最小有效组合。LoRA、Cross-Attention 等未被保留的方案并非作为失败记录，而是作为研究路线收敛的重要证据。

本文通过逐步筛选的方式回答以下基础问题：首先，性能提升究竟主要来自训练信号优化还是结构改动；其次，在训练信号改善之后，剩余误差更接近全局语义不足还是局部边界不足；最后，在可选的多条增强路线中，哪些方案真正与当前任务匹配。

\section{研究目标与技术路线}

\begin{figure}[htbp]
  \centering
  \includegraphics[width=0.85\textwidth]{figures/tech_roadmap.pdf}
  \caption{本文技术路线总览图。图中各方案标注的数值为概要性展示，具体实验结果以正文表格为准。}
  \label{fig:tech_roadmap}
\end{figure}

如图~\ref{fig:tech_roadmap}所示，本文按照”损失策略筛选—适配路线筛选—最终结果验证”的论证顺序展开，而不是简单按实验发生顺序罗列结果。全部实验均在 GT box prompt 条件下进行统一评估，并以 15 个测试病例（10 个 FLARE22 + 5 个 AMOS22）为病例级统计对象，以保证不同方案之间的内部可比性。具体研究步骤如下：

\begin{enumerate}
  \item \textbf{损失策略筛选}：围绕训练信号失衡问题，设计 Balance Loss，并通过组件消融、控制变量验证和补充对照实验，筛选后续统一采用的损失主干。
  \item \textbf{适配路线筛选}：在固定最优损失配置后，比较 Cross-Attention、LoRA 不同秩设置和局部适配器等路线，回答“何种结构增强方式最适合当前任务”。
  \item \textbf{最终结果验证}：仅针对最终保留的里程碑方案，展示总体结果、逐器官表现、边界质量、定性可视化和局限性分析，形成完整的证据闭环。
\end{enumerate}

该技术路线具有明确的问题导向：第一阶段确认训练信号是否为当前性能瓶颈的主要来源，第二阶段判断在训练信号重构后剩余误差是否主要表现为局部边界表征不足。此外，本文保留 BL-$\alpha$1.0 退化、Cross-Attention 回落、冻结主干 LoRA 效果受限等负向结果作为路线收敛的排除性证据，以增强最终方案的说服力。

\section{主要工作与贡献}

围绕上述目标，本文的主要工作与贡献可概括为以下三点：

\begin{enumerate}
  \item \textbf{设计并系统验证了面向训练信号重构的 Balance Loss。} 针对前景/背景失衡和困难像素学习不充分问题，本文设计了将类别间重平衡（Inter-CBL）与类别内困难像素加权（Intra-CBL）联合建模的 Balance Loss，并辅以两阶段训练策略控制引入时机。通过系统的组件消融、超参数敏感性分析和竞争对照实验，确定 BL 为后续统一采用的基础损失配置，为后续结构路线比较提供稳定基座。
  \item \textbf{完成了面向 MedSAM 的多条微调/适配路线的系统比较，并保留局部适配方案。} 本文将 LoRA、Cross-Attention 融合与局部适配器等路线统一纳入比较框架，而非预设某一路线为最优。实验表明，在当前任务和配置下，局部适配方案更能补足边界相关表征，从而形成”Balance Loss + 局部适配”的最终组合。LoRA 和 Cross-Attention 的负向结果作为路线收敛的排除性证据予以保留。
  \item \textbf{通过系统消融与路线对比，建立了当前任务中不同优化策略的适用性判据。} 实验表明，训练信号重构是整体性能提升的首要来源，而局部边界补强在此基础上提供有效补充；Cross-Attention 融合与冻结主干 LoRA 微调等方案的负向结果进一步明确了不同策略在当前任务中的适用边界，使最终方案的选择具有更强的可解释性。
\end{enumerate}

本文的贡献既体现在 Balance Loss 与局部适配器两条方法主线上，也体现在对负向结果的保留与不同策略适用边界的澄清上。后续章节将 BL 与 BL+MSL 作为阶段性里程碑方案重点展开。

\section{论文组织结构}

全文共分六章：

\begin{itemize}
  \item 第1章为绪论，介绍研究背景、相关工作、论文定位与整体技术路线。
  \item 第2章介绍相关理论基础，包括 MedSAM 架构原理、不平衡学习理论、PEFT 技术及局部适配相关理论。
  \item 第3章聚焦 Balance Loss 的配置选择，回答“损失策略如何被筛选出来”。
  \item 第4章聚焦微调与结构增强路线比较，回答“为什么最终选择局部适配而不是其他路线”。
  \item 第5章展示最终保留方案的整体结果、逐器官表现、边界分析、定性可视化与局限性讨论。
  \item 第6章总结全文工作，并讨论后续可进一步扩展的研究方向。
\end{itemize}

\section{本章小结}

本章围绕本文的研究问题、研究动机与论证组织给出了整体说明。本文以 MedSAM 为基座，以腹部多器官 CT 为统一验证场景，围绕训练信号重构与局部边界补强两条主线展开优化，并将实验章节组织为先筛选后展示的研究型叙事结构，为后续章节的展开奠定基础。
